% Delta de §10.27. Insertar al final de 05_ciclo_catalan.tex.
% La actualización afín y su balance permanecen en sus propietarios originales.
\subsection{Traces de retour et moments du quotient constitutif}
\label{amp-afin-sec}

La mise à jour du registre et la réponse du cycle partagent le
transport dodécaphasique et conservent des informations de types différents.
L’événement signé s’obtient à partir de la direction et de l’orientation de chaque
fenêtre de l’histoire APP--TRIT--TPK. Son accroissement vectoriel et le
bilan~\eqref{mem:balance} retiennent l’interaction avec le registre préalable.
La construction suivante prolonge la lecture spectrale du même cycle,
une fois ses sous-espaces d’incidence fixés. Cet ordre permet de composer
le retour, le déterminant et les moments en conservant l’histoire qui les précède.

Nous conservons l’opérateur $C$ et les projecteurs de
\eqref{iii:eq:projectors34}, et définissons
\begin{equation}
 E_{\rm cyc}=P_3^{\mathrm{cyc}}-P_4^{\mathrm{cyc}},\qquad
 b_n=\operatorname{Tr}(C^n E_{\rm cyc}),\qquad n\geq0.
 \label{amp-afin-trazas}
\end{equation}
Le symbole $b_n$ distingue cette trace spectrale de l’événement signé
$a_m$ d’une fenêtre. La trace utilise l’opérateur de transport ;
l’événement utilise en outre la trajectoire effectivement parcourue.

\begin{proposition}[Tableau des traces produit par l’incidence]
\label{amp-afin-prop-tabla}
On a
\begin{equation}
 \begin{gathered}
 b_n=3\mathbf1_{3\mid n}-4\mathbf1_{4\mid n},\\
 (b_1,\ldots,b_{12})=(0,0,3,-4,0,3,0,-4,3,0,0,-1).
 \end{gathered}
 \label{amp-afin-tabla}
\end{equation}
La période minimale est douze et la somme sur une période est nulle. Pour
$B_N=\sum_{n=1}^{N}b_n$,
\begin{equation}
 B_N=3\lfloor N/3\rfloor-4\lfloor N/4\rfloor,\qquad |B_N|\leq3.
 \label{amp-afin-parciales}
\end{equation}
\end{proposition}
\begin{proof}
Dans $H_d$, l’opérateur $C$ permute cycliquement les $d$ classes de
coordonnées. Sa puissance $n$ a $d$ points fixes si $d$ divise $n$
et aucun dans le cas contraire. Comme $P_d^{\mathrm{cyc}}$ commute avec $C$,
$\operatorname{Tr}(C^nP_d^{\mathrm{cyc}})$ est la trace de cette restriction.
Soustraire les cas $d=3,4$ prouve la formule. La période positive minimale
est douze : $b_n=-1$ se produit exactement aux multiples de douze.
Sur une période, les contributions sont $3\cdot4-4\cdot3=0$.
La somme finie compte les multiples et donne~\eqref{amp-afin-parciales}.
Ses valeurs pour $N=0,\ldots,11$ sont
$(0,0,0,3,-1,-1,2,2,-2,1,1,1)$ et se répètent lorsqu’on ajoute douze.
\end{proof}

\begin{theorem}[Reconstruction logarithmique du lecteur constitutif]
\label{amp-afin-thm-logdet}
Pour $|s|<1$, le quotient $f$ de~\eqref{iii:eq:trace-det} vérifie
\begin{align}
 -\log f(s)&=\sum_{n\geq1}\frac{b_n}{n}s^n,\label{amp-afin-logdet}\\
 -\frac{s f'(s)}{f(s)}
 &=\frac{3s^3}{1-s^3}-\frac{4s^4}{1-s^4}.
 \label{amp-afin-logder}
\end{align}
La branche holomorphe est fixée par $\log f(0)=0$.
Sur $0<s<1$, elle coïncide avec le logarithme réel.
\end{theorem}
\begin{proof}
Les polynômes $1-s^3$ et $1-s^4$ ne s’annulent pas dans le disque ouvert.
Le développement $-\log(1-w)=\sum_{k\geq1}w^k/k$, obtenu en intégrant
la série géométrique depuis zéro, converge uniformément sur les disques
compacts. Avec $w=s^3,s^4$, leur différence donne
$-\log f(s)=\sum_{k\geq1}s^{3k}/k-\sum_{k\geq1}s^{4k}/k$.
Le coefficient de $s^n$ est $b_n/n$.
La convergence locale uniforme permet de dériver et de sommer les deux séries
géométriques, ce qui prouve~\eqref{amp-afin-logder}.
La condition en zéro fixe la constante d’intégration et la branche.
\end{proof}

Ainsi, $b_0/12=-1/12$ enregistre le défaut normalisé de dimension,
tandis que toute la suite $b_n$ reconstruit une fonction.
Ce sont deux lectures de portées différentes de la même paire de sous-espaces.
Lorsqu’on substitue $s=q_\pm^{30}$, les arguments restent les sorties
de la construction angulaire préalable ; le développement ne les sélectionne pas.

\begin{theorem}[Moments de profondeur et représentation de Mellin]
\label{amp-afin-thm-momentos}
La série
\begin{equation}
 M_b(z)=\sum_{n\geq1}\frac{b_n}{n^z}
 \label{amp-afin-momento}
\end{equation}
converge localement uniformément pour $\operatorname{Re}z>0$.
Si $\sigma=\operatorname{Re}z>0$, son reste satisfait
\begin{equation}
 \left|M_b(z)-\sum_{n=1}^{N}\frac{b_n}{n^z}\right|
 \leq3\left(1+\frac{|z|}{\sigma}\right)(N+1)^{-\sigma}.
 \label{amp-afin-cota}
\end{equation}
Dans le domaine de convergence absolue $\operatorname{Re}z>1$,
\begin{equation}
 M_b(z)=(3^{1-z}-4^{1-z})\zeta(z),\qquad
 \zeta(z)=\sum_{n\geq1}n^{-z}.
 \label{amp-afin-reconocimiento}
\end{equation}
Pour $\operatorname{Re}z>0$, on a en outre
\begin{equation}
 \Gamma(z)M_b(z)=
 \int_0^\infty t^{z-1}
 \left(\frac3{e^{3t}-1}-\frac4{e^{4t}-1}\right)dt,
 \label{amp-afin-mellin}
\end{equation}
où $\Gamma(z)=\int_0^\infty u^{z-1}e^{-u}\,du$.
\end{theorem}
\begin{proof}
La sommation par parties et $|B_n|\leq3$ donnent, en passant à la limite supérieure,
\[
 \sum_{n>N}b_n n^{-z}
 =-B_N(N+1)^{-z}
   +\sum_{n>N}B_n\bigl(n^{-z}-(n+1)^{-z}\bigr).
\]
La différence de puissances est bornée en module par
$|z|\int_n^{n+1}x^{-\sigma-1}dx$. Cela prouve
\eqref{amp-afin-cota}, la convergence et son uniformité sur les compacts.
Pour $\sigma>1$, la convergence absolue permet de regrouper
les multiples de trois et de quatre séparément, ce qui donne
\eqref{amp-afin-reconocimiento}.

La série géométrique du transport, évaluée en $s=e^{-t}$, produit
\[
 \sum_{n\geq1}b_n e^{-nt}
 =\frac3{e^{3t}-1}-\frac4{e^{4t}-1}=:K_b(t).
\]
Pour $\sigma>1$, on peut intégrer terme à terme parce que
$\sum|b_n|n^{-\sigma}<\infty$. La substitution $u=nt$ donne
$\Gamma(z)n^{-z}$ et prouve~\eqref{amp-afin-mellin} dans ce demi-plan.
À l’origine, les termes $1/t$ s’annulent :
$K_b(t)=1/2+O(t)$. À l’infini, $K_b(t)=O(e^{-3t})$.
L’intégrale est donc holomorphe pour $\operatorname{Re}z>0$
par convergence locale uniforme. La série l’est aussi par la borne
précédente ; le principe d’identité prolonge l’égalité à ce domaine.
\end{proof}

En $0<\operatorname{Re}z\leq1$, l’ordre croissant de profondeur
fixe la lecture convergente de la série. La réalisation comme trace
absolue de l’opérateur diagonal correspond à $\operatorname{Re}z>1$.

\begin{corollary}[Évaluations compatibles avec l’extrémité constitutive]
\label{amp-afin-cor-extremos}
On a
\begin{equation}
 M_b(2)=\frac{\zeta(2)}{12},\qquad
 M_b(1)=\log(4/3)=-\log f(1),\qquad f(1)=\frac34.
 \label{amp-afin-extremos}
\end{equation}
\end{corollary}
\begin{proof}
La première égalité résulte de $3^{-1}-4^{-1}=1/12$.
Pour la seconde, les sommes partielles exactes sont
$H_{\lfloor N/3\rfloor}-H_{\lfloor N/4\rfloor}$, avec
$H_j=\sum_{k=1}^j1/k$ et $H_0=0$.
La comparaison intégrale de $1/x$ entre ces deux indices donne la limite
$\log(4/3)$. Annuler $1-s$ avant d’évaluer le quotient déterminantal
donne $f(1)=3/4$.
\end{proof}

Le registre affine retient les événements et leur interaction avec la mémoire
préalable ; le tableau des traces retient l’incidence des deux sous-espaces
à toutes les puissances du transport ; le moment applique une pondération
de profondeur à ce tableau déjà produit. Sa périodicité spectrale
coexiste avec une histoire d’événements qui se poursuit après le retour
de phase. Inverser $C$ conserve $b_n$, parce que $d\mid n$ est invariant
sous $n\mapsto-n$ ; le coefficient orienté de la résolvante de
\eqref{iii:eq:resolvent}, avec ses marques fixes, change de signe.
La lecture conjointe conserve une distinction que le déterminant
scalaire n’enregistre pas à lui seul. La forme quadratique et la somme des
coordonnées du registre maintiennent leurs types mathématiques ; cette
composition ne leur assigne pas d’unités d’énergie ni de charge électrique.

\subsection{Composition résiduelle et produit de moments}
\label{amp-afin-sec-productos}

La composition ultérieure des lecteurs conserve également l’opération
effectuée sur la profondeur. Dans la réalisation
$\mathcal H=\ell^2(\mathbb N_{>0})$, nous fixons
$N\delta_n=n\delta_n$. Soient $a,c$ deux tableaux résiduels périodiques
déjà produits et $Q_a\delta_n=a(n)\delta_n$,
$Q_c\delta_n=c(n)\delta_n$. Lorsqu’on les prolonge tous deux à la période commune,
les deux lecteurs agissent sur le même indice. Pour
$\operatorname{Re}z>1$, nous écrivons
$M_a(z)=\operatorname{Tr}(N^{-z}Q_a)$ et de même $M_c$.

\begin{proposition}[Composition diagonale et composition tensorielle]
\label{amp-afin-prop-productos}
On a les identités
\begin{align}
 \operatorname{Tr}(N^{-z}Q_aQ_c)
 &=\sum_{n\geq1}\frac{a(n)c(n)}{n^z},
 \label{amp-afin-diagonal}\\
 M_a(z)M_c(z)
 &=\operatorname{Tr}_{\mathcal H\otimes\mathcal H}
 \bigl((N\otimes N)^{-z}(Q_a\otimes Q_c)\bigr)\notag\\
 &=\sum_{k\geq1}\frac{(a*c)(k)}{k^z},
 \label{amp-afin-tensor}\\
 (a*c)(k)&=\sum_{d\mid k}a(d)c(k/d),
 \qquad \operatorname{Re}z>1.
 \label{amp-afin-convolucion}
\end{align}
Tous les opérateurs figurant à l’intérieur de ces traces sont à trace.
\end{proposition}
\begin{proof}
Le produit diagonal agit sur $\delta_n$ au moyen de
$a(n)c(n)$, ce qui prouve~\eqref{amp-afin-diagonal}.
Dans $\delta_m\otimes\delta_n$, l’opérateur $N\otimes N$
a pour valeur propre $mn$ et $Q_a\otimes Q_c$ agit au moyen de
$a(m)c(n)$. Si $\sigma=\operatorname{Re}z>1$, la somme des
modules de ces coefficients spectraux est bornée par
$\|a\|_\infty\|c\|_\infty(\sum n^{-\sigma})^2<\infty$.
Cette borne prouve la condition de classe de trace et permet de factoriser
la somme double ou de la regrouper par $k=mn$. Les paires de produit $k$
sont exactement $(d,k/d)$, avec $d\mid k$.
La borne analogue avec une seule somme prouve le cas diagonal.
\end{proof}

Dans les lecteurs de cette section, prendre $a(n)=b_n$ et
$c(n)=\langle v,C^nu\rangle$ donne, pour $n=3$,
$a(3)c(3)=-3$, tandis que $(a*c)(3)=3$.
Cette différence de coefficients met en évidence l’information conservée par
chaque opération. La composition diagonale maintient une profondeur commune ;
la tensorielle conserve deux profondeurs avant de les regrouper par leur produit.
La réalisation tensorielle n’attribue pas d’indépendance statistique aux
histoires HMT qui ont produit les lecteurs. Les dérivées régularisées et
les changements de normalisation d’autres moments constituent des opérations
ultérieures distinctes des deux compositions démontrées ici.
